\chapter{Sept.~18 --- Inversion Formula}

\section{The Fej\'er Kernel}

\begin{definition}
  Given $f \in L^1(\R)$, the
  \emph{inverse Fourier transform} of $f$ is
  \[
    \widecheck{f}(\xi)
    = \int f(x) e^{2\pi i \xi x}\, dx
    = \widehat{f}(-\xi), \quad \xi \in \R.
  \]
\end{definition}

\begin{remark}
  Recall the Riemann-Lebesgue lemma says that
  if $f \in L^1(\R)$, then
  $\widehat{f} \in C_0(\R)$. However, we need not in
  general have $\widehat{f} \in L^1(\R)$.
  For example,
  \[
    \widehat{\chi}_{[-T, T]}(\xi)
    = \frac{\sin 2\pi T \xi}{\pi \xi}
    = d_{2\pi T}(\xi) \notin L^1(\R),
  \]
  where $d(\xi) = (\sin \pi \xi) / \pi \xi$
  is the Dirichlet function. In particular,
  we cannot compute $(d_{2\pi T})^\wedge$
  or $(d_{2\pi T})^\vee$ now since $d_{2\pi T}$
  is not integrable. However, $d_{2\pi T} \in L^2(\R)$:
  Later we will extend the
  Fourier transform and its inverse to $L^2(\R)$, and
  we will see that in that context,
  \[
    (d_{2\pi T})^{\wedge} = \chi_{[-T, T]} = (d_{2\pi T})^\vee.
  \]
  Then we will see that in the $L^2(\R)$ sense,
  $(f * d_{2\pi T})^\vee = \widehat{f} \cdot (d_{2\pi T})^\wedge = \widehat{f} \chi_{[-T, T]}$.
\end{remark}

\begin{definition}
  The \emph{Fej\'er function} is
  $w(x) = (\sin (\pi x) / \pi x)^2 = (d_{\pi} (x))^2$.
  The \emph{Fej\'er kernel} is
  \[w_\lambda(x) = \lambda w(\lambda x).\]
\end{definition}

\begin{remark}
  Note that $w \in L^1(\R)$, and so
  $\{w_\lambda\}_\lambda$ is an approximate identity.
  One can check that
  \[
    \int_{-\infty}^\infty w(x)\, dx = 1.
  \]
  Let $\chi = \chi_{[-1 / 2, 1 / 2]} \in L^1(\R)$
  with $\int_{-\infty}^\infty \chi = 1$.
  Then $\widehat{\chi}(\xi) = d_\pi(\xi) = \sin(\pi \xi) / \pi \xi$.
  Set
  $W = \chi * \chi$, with
  \[
    W(\xi) = \max\{1 - |\xi|, 0\}
  \]
  (the graph of $W$ looks like a triangle supported
  on $[-1, 1]$). We then have
  \[
    \widehat{W}(\xi)
    = (\chi * \chi)^\wedge(\xi)
    = (\widehat{\chi}(\xi))^2
    = (d\pi(\xi))^2 = w(\xi).
  \]
  Note that if we knew the inversion formula,
  then we would get
  $\widehat{w} = (\widecheck{W})^\wedge = W$ also.
\end{remark}

\section{Inversion Formula}

\begin{remark}[Motivation for proof strategy]
  Fix $f \in L^1(\R)$. Since $\{w_\lambda\}$
  is an approximate identity,
  \[
    f * w_\lambda \longrightarrow f
  \]
  in $L^1$-norm. If we knew that
  $\widehat{w} = W$, then on the Fourier side,
  \[
    (f * w_\lambda)^\wedge(\xi)
    = \widehat{f}(\xi) \widehat{w}_\lambda(\xi)
    = \widehat{f}(\xi) \widehat{w}(\xi / \lambda)
    = \widehat{f}(\xi) W(\xi / \lambda),
  \]
  where $W(\xi / \lambda) \to 1$ pointwise as
  $\lambda \to \infty$.
  So if we knew the inversion formula, then
  \[
    (f * w_\lambda)(x)
    = \int_{-\infty}^\infty \widehat{f}(\xi) W(\xi / \lambda) e^{2\pi i \xi x}\, d\xi,
  \]
  and taking $\lambda \to \infty$ we get the conclusion.
  We will instead proceed to show the above identity
  directly.
\end{remark}

\begin{theorem}
  If $f \in L^1(\R)$, then
  \[
    (f * w_\lambda)(x)
    = \int_{-\infty}^\infty \widehat{f}(\xi) W(\xi / \lambda) e^{2\pi i \xi x}\, dx
    = \int_{-\lambda}^\lambda \widehat{f}(\xi) \left(1 - \frac{|\xi|}{\lambda}\right) e^{2\pi i \xi x}\, d\xi.
  \]
\end{theorem}

\begin{proof}
  By definition,
  \[
    (f * w_\lambda)(x) = \int_{-\infty}^\infty f(y) w_\lambda(x - y)\, dy.
  \]
  We know that $w(x) = \widehat{W}(x) = \widecheck{W}(x)$,
  and so
  \[
    w_\lambda(x)
    = \lambda w(\lambda x)
    = \lambda \widecheck{W}(\lambda x)
    = \lambda \int_{-\infty}^\infty W(\xi) e^{2\pi i \xi \lambda x}\, d\xi
    = \int_{-\infty}^\infty W(\xi / \lambda) e^{2\pi i \xi x}\, d\xi
  \]
  after a change of variables (replacing
  $\lambda x$ with $x$). So
  \[
    (f * w_\lambda)(x)
    = \int_{-\infty}^\infty f(y) \int_{-\infty}^\infty W(\xi / \lambda) e^{2\pi i \xi (x - y)}\, d\xi dy
    = \int_{-\infty}^\infty \int_{-\infty}^\infty f(y) W(\xi / \lambda) e^{2\pi i \xi (x - y)}\, d\xi dy.
  \]
  Since the integrand is integrable in $\xi, y$, by
  Fubini's theorem
  \[
    (f * w_\lambda)(x)
    = \int_{-\infty}^\infty \left(\int_{-\infty}^\infty f(y) e^{-2\pi iy}\, dy\right) W(\xi / \lambda) e^{2\pi i \xi x}\, d\xi
    = \int_{-\infty}^\infty \widehat{f}(\xi) W(\xi / \lambda) e^{2\pi i \xi x}\, d\xi.
  \]
  The last equality in the desired result follows from
  $W(\xi / \lambda) = (1 - |\xi| / \lambda) \chi_{[-\lambda, \lambda]}$.
\end{proof}

\begin{theorem}[Inversion formula]
  If $f, \widehat{f} \in L^1(\R)$, then
  $f, \widehat{f} \in C_0(\R)$ and
  \[
    f(x) = (\widehat{f})^\vee(x) = \int_{-\infty}^\infty \widehat{f}(\xi) e^{2\pi i \xi x}\, d\xi
    \quad \text{and} \quad
    f(x) = (\widecheck{f})^\wedge(x)
    = \int_{-\infty}^\infty \widecheck{f}(\xi) e^{2\pi i \xi x}\, d\xi,
  \]
  with both equalities holding pointwise at
  every $x \in \R$.
\end{theorem}

\begin{proof}
  We already have
  \[
    (f * w_\lambda)(x)
    = \int_{-\infty}^\infty \widehat{f}(\xi) W(\xi / \lambda) e^{2\pi i \xi x}\, d\xi.
  \]
  The integrand on the right-hand side
  converges pointwise
  to $\widehat{f}(\xi) e^{2\pi i \xi x}$
  as $\lambda \to \infty$
  (since $W(\xi / \lambda) \to 1$).
  Since
  $|\widehat{f}(\xi) W(\xi / \lambda) e^{2\pi i \xi x}| \le |\widehat{f}(\xi)|$
  and $\widehat{f} \in L^1(\R)$,
  the dominated convergence theorem gives
  \pagebreak
  \[
    (f * w_\lambda)(x)
    \xrightarrow[\lambda \to \infty]{}
    \int_{-\infty}^\infty \widehat{f}(\xi) e^{2\pi i \xi x}\, d\xi
  \]
  pointwise for every $x \in \R$ (recall that
  $f \in L^1(\R)$ and $w_\lambda \in C_c(\R)$ implies
  $f * w_\lambda \in C_0(\R)$). Also, since
  $\{w_\lambda\}$ is an approximate identity,
  $f * w_\lambda \to f$ in $L^1$-norm as
  $\lambda \to \infty$. Then there exists a
  subsequence $f * w_{\lambda_n} \to f$ pointwise
  a.e., so by the uniqueness of limits
  \[
    f(x) = \int_{-\infty}^\infty \widehat{f}(\xi) e^{2\pi i \xi x}\, d\xi
    = (\widehat{f})^\vee(x)
  \]
  for a.e. $x \in \R$. Since $\widehat{f} \in L^1(\R)$,
  we have
  $(\widehat{f})^\vee \in C_0(\R)$, and so by
  redefining $f$ on a set of measure zero
  we may arrange for the above equality to hold for every
  $x \in \R$.
\end{proof}

\begin{corollary}
  We have
  $\widehat{w} = (\widecheck{W})^\wedge = W = (\widehat{W})^\vee = \widecheck{w}$.
\end{corollary}

\begin{proof}
  This follows from the inversion formula,
  since $w, W$ are integrable and
  we know $\widehat{W} = w$.
\end{proof}

\begin{corollary}[Uniqueness theorem]
  If $f, g \in L^1(\R)$, then
  $f = g$ a.e. if and only if $\widehat{f} = \widehat{g}$
  a.e. In particular,
  $f = 0$ a.e. if and only if $\widehat{f} = 0$ a.e.
\end{corollary}

\begin{proof}
  The first statement follows from the latter
  by linearity, and
  clearly $\widehat{f} = 0$ a.e. if $f = 0$ a.e.
  Now let $f \in L^1(\R)$ with
  $\widehat{f} = 0$ a.e. Then $\widehat{f} = L^1(\R)$,
  so by the inversion formula,
  $f = (\widehat{f})^\vee = \widecheck{0} = 0$ a.e.
\end{proof}

\begin{remark}
  If $\widehat{f} = 0$ a.e., then we could also directly
  see from
  \[
    (f * w_\lambda)(x)
    = \int_{-\infty}^\infty \widehat{f}(\xi) W(\xi / \lambda) e^{2\pi i \xi x}\, d\xi
    = 0
  \]
  that $f = 0$ a.e. since $f * w_\lambda \to f$
  in $L^1$-norm.
\end{remark}
